\chapter{Atomic Nuclei Can Change Too}

\begin{kaichang}
Many homes have a small white disk near a ceiling corner: a smoke alarm. Turn it over and you may find tiny print saying it contains a trace of the radioactive substance americium-241. Radioactive! The word conjures radiation symbols, protective suits, and sirens, yet this substance lives openly outside your bedroom?

Now look at the banana in your hand. Bananas are rich in potassium, and natural potassium contains a tiny amount of potassium-40, also radioactive. Right now, nuclei in your hand and body are quietly transforming.

Do not panic or guess. By this chapter's end, you will understand why a smoke alarm needs a decaying nucleus and whether that banana should worry you.
\end{kaichang}

\section{An Invisible Leak}

In Paris in 1896, the physicist Becquerel locked uranium-bearing minerals and photographic plates tightly wrapped in black paper in a drawer. He intended to study fluorescence, but poor weather postponed the experiment. A few days later, he developed the plates anyway and found clear mineral outlines. They had been exposed despite never seeing any light.

Humanity had encountered a leak that could not be seen, heard, or smelled. Uranium minerals continuously emitted something that penetrated black paper and exposed photographic plates. Later, the Curies found uranium ore residues that leaked more intensely than pure uranium. Processing tons of material, they isolated two new elements, polonium and radium, and named the phenomenon radioactivity.

Today you can hear it without touching a mineral. Bring a Geiger counter, an instrument for detecting radiation, near an old radium watch dial or certain mineral specimens, and the speaker crackles like rapid rain on sheet metal. Each click catches a nucleus transforming at that instant.

Safety must come first. We will say this once, and every word matters: \textbf{all radioactive material discussed here is for reading, videos, or teacher demonstrations only. Never attempt these activities at home.} Do not dismantle smoke alarms, collect unidentified luminous antiques or minerals, or buy anything advertised as a radiation source online. Radiation is invisible and odorless, and you have no built-in alarm. That is precisely its danger, and why professional equipment and personnel must handle it.

The good news is that you are about to learn something more useful than an alarm: a set of tools for distinguishing real danger from needless fear.

\section{A Tug-of-War Inside the Nucleus}

Why do some nuclei leak while others remain unchanged for trillions of years? The answer lies in a ceaseless nuclear tug-of-war.

Recall Chapter 6's picture: a nucleus is only one hundred-thousandth of an atom's size yet holds nearly all its mass. Inside are positive protons and uncharged neutrons. Trouble is immediate: like charges repel, so a crowd of protons in such tiny space experiences electromagnetic forces trying desperately to push them apart. The nucleus stays together because at sufficiently close range another, much stronger force operates: the strong interaction, binding protons and neutrons like superglue.

The crucial point is that these forces have entirely different personalities:

\begin{itemize}[itemsep=2pt]
\item Electromagnetic repulsion reaches far. Every nuclear proton repels every other, so larger nuclei accumulate more repulsion.
\item The strong interaction acts only between close neighbors. Its range is extremely short; beyond arm's reach it has no influence.
\end{itemize}

A nucleus's fate depends on their score. In a small nucleus, few, close neighbors let glue dominate, giving stability. As the nucleus grows, repulsion's total rises faster because it adds across all protons, while glue acts only locally. Beyond a threshold, even powerful glue cannot hold order. The nucleus becomes unstable and may suddenly eject a fragment or rearrange its internal personnel, becoming a lower-energy, more stable nucleus. This is radioactive decay: \textbf{the spontaneous restructuring of an unstable nucleus accompanied by radiation}.

Notice the connection to Chapter 7. One element, with a fixed proton count, may have several isotopes distinguished by neutrons. Neutrons contribute glue without electromagnetic repulsion, so their proper proportion directly affects stability. Carbon-12, with 6 protons and 6 neutrons, is rock-steady, but carbon-14, with 6 and 8, is radioactive. One elemental identity card, two very different nuclear fates.

\section{Three Decays, Three Different Accounts}

Unstable nuclei mainly transform by three routes, corresponding to three radiations. Their bookkeeping differs completely and deserves careful examination.

\textbf{$\alpha$ decay: ejecting a helium nucleus.} An oversized nucleus often throws out a package of 2 protons and 2 neutrons, exactly a helium nucleus. The parent loses 4 from its mass number and 2 from its proton count. Since proton count is elemental identity, \textbf{the element changes}. A smoke alarm's americium-241 follows this route:
\[{}^{241}_{95}\chem{Am} \rightarrow {}^{237}_{93}\chem{Np} + {}^{4}_{2}\chem{He}\]
Americium becomes neptunium. Large, highly charged $\alpha$ particles do not travel far: paper, a few centimeters of air, or outer skin can stop them. But if their source is swallowed or inhaled, close-range damage can be severe.

\textbf{$\beta$ decay: a neutron becomes a proton.} In a neutron-rich nucleus, one neutron can become a proton while emitting a fast electron, the $\beta$ particle, and an almost undetectable antineutrino. Mass number stays fixed while proton number rises by 1: \textbf{the element changes again}, moving one place along. Archaeologists' carbon-14 becomes nitrogen this way:
\[{}^{14}_{6}\chem{C} \rightarrow {}^{14}_{7}\chem{N} + {}^{0}_{-1}\chem{e}\]
Much smaller than $\alpha$ particles, $\beta$ particles penetrate outer skin but can be stopped by a few millimeters of aluminum.

\textbf{$\gamma$ decay: releasing energy without changing composition.} After either preceding transformation, a new nucleus often remains excited at high energy. It calms by emitting extremely energetic electromagnetic radiation, $\gamma$ rays. Proton and neutron counts stay fixed, so the element remains unchanged. But $\gamma$ rays penetrate most strongly, requiring thick lead or concrete for substantial attenuation. Much medical and industrial radiation is of this kind.

Together, these routes reveal something that would have astonished ancient thinkers: \textbf{nuclei change, and elements transform into other elements}. Medieval alchemists dreamed of turning lead into gold but could not budge a nucleus in a lifetime, because chemistry changes external electrons and leaves nuclei untouched. Decay genuinely transmutes elements: uranium undergoes a series of transformations and ultimately becomes stable lead. The direction and rate, however, are not ours to choose. Note also that stars use another route, fusion, to forge heavier elements from lighter ones. Your carbon and iron came from this process, the story of Chapter 16.

\section{Half-Life: About Populations, Not Individuals}

The next question follows naturally: when will a particular carbon-14 nucleus decay?

The answer is \textbf{nobody knows, and there is no way to know}. This is not an instrumental limitation. Nuclear decay is fundamentally random: a nucleus destined to decay tomorrow and one a century from now are indistinguishable today. Watching cannot hurry it; freezing, heating, electrifying, and compressing it do not help.

Yet nature leaves a window: \textbf{individuals are unpredictable, but populations behave with extraordinary precision}. In a large collection of identical radioactive nuclei, the remaining number halves after each fixed interval, the half-life. Carbon-14's is about 5730 years: one gram becomes half a gram after 5730 years, a quarter after 11460, then an eighth, a sixteenth, and so on. Iodine-131's half-life is only 8 days, polonium-214's 0.00016 seconds, while uranium-238's is 4.5 billion years, almost Earth's age.

\begin{qiaoliang}
Why does the amount halve in equal intervals? Each nucleus has a constant probability of decaying per unit time. Think of coins: each has a one-half chance of heads, so one toss eliminates about half a handful; toss the survivors and another half goes. Nuclear decay is like each nucleus continuously tossing such a coin. With $N_0$ nuclei initially, elapsed time $t$, and half-life $T$, the remainder is
\[N = N_0 \times \left(\frac{1}{2}\right)^{t/T}\]
Take iodine-131, with $T=8$ days. After 80 days of hospital storage, $t/T=10$, leaving $(1/2)^{10}=1/1024$, less than a thousandth. This is why hospitals are willing to treat patients with it: after a few weeks, almost nothing remains in the body.
\end{qiaoliang}

Estimate how large the population really is. Living organisms exchange carbon with their surroundings, and each gram contains about fifty billion carbon-14 atoms. A 5730-year half-life is about $1.8\times 10^{11}$ seconds, giving a per-second decay fraction of approximately $\ln 2 \div (1.8\times 10^{11}) \approx 4\times 10^{-12}$. Multiply by fifty billion: each gram of your body's carbon undergoes about 14 carbon-14 decays per minute. \textbf{You yourself are a faintly clicking Geiger counter}, only the activity is tiny and completely harmless.

\begin{shiyan}
\textbf{Model Half-Life with a Handful of Coins} (Safe to do on a table.)

Materials: 50--100 identical coins, buttons, or Go stones, a box, and paper.

Procedure: Put all the coins in the box, shake, and pour them out. Remove every heads-up coin: those have decayed. Record the decayed and remaining counts. Return the survivors to the box and repeat until fewer than ten remain. Each round represents one half-life.

What to observe: (1) Each round removes roughly half the remainder, but not exactly half. Plot remaining number against round number and see whether it resembles a slide that repeatedly halves. (2) Repeat with only 4 coins. Is the curve still smooth? Notice that precision emerges only in large groups. (3) Choose one coin: can you predict its elimination round?

Safety: This experiment uses no radioactivity. Observe real radiation, such as Geiger counting or cloud-chamber particle tracks, only through teacher demonstrations or reputable science videos.
\end{shiyan}

\begin{zhengju}
``An element spontaneously becomes another element.'' When Rutherford and his assistant Soddy proposed this in 1902, it nearly amounted to tearing up chemistry's foundations in public. Since Dalton, the consensus had rested on indivisible atoms and unchanging elements.

The clues were a chain of odd observations. After the Curies discovered radium, people noticed a radioactive gas appearing around radioactive substances, later named radon. Radon itself gradually died away, leaving a new radioactive film on vessel walls. Rutherford and Soddy traced this chain of radioactivity producing radioactivity and measured each activity's rise and fall. They concluded that substances were not catching activity like an infection: nuclei were transforming step by step, each becoming a new element and emitting radiation. Uranium became thorium, thorium protactinium, protactinium another uranium isotope, continuing until lead.

Opposition was fierce because the explanation violated two cherished beliefs. But the test was straightforward: if decay made new elements, its endpoint, lead, should accumulate with the age of uranium ores. Helium, formed when $\alpha$ particles acquire electrons, should accumulate too. Geologists examined ancient ores and found lead in amounts matching their ages and abundant helium bubbles. Prediction fulfilled, transmutation moved from heresy to textbook. The point is not who was clever: however radical a conclusion sounds, specific testable predictions let science decide.
\end{zhengju}

\section{Where Energy Comes From: Fission, Fusion, and Binding Energy}

Decay radiation carries substantial energy, but nuclei hold a much greater treasure. To understand it, remember: \textbf{separating a nucleus requires energy; assembling dispersed particles into one releases energy}. The more favorable the assembly, the more is released.

Nuclei differ in how favorable their assembly is. Plot binding energy per nucleon, the average strength holding each proton or neutron, and you see a mountain high in the middle and low at both ends. Iron stands near the summit, bound most tightly of all, the energy valley bottom. This points to two directions for energy release:

\begin{itemize}[itemsep=2pt]
\item \textbf{Fission}: nuclei much heavier than iron, such as uranium-235, split into two medium-sized fragments. This moves toward the valley bottom: nucleons bind more tightly and release the difference as energy. A neutron striking uranium-235 can trigger fission, which ejects 2--3 more neutrons to strike other nuclei. The chain continues by itself: a chain reaction. A nuclear power plant controls its pace precisely with control rods, letting heat emerge steadily to boil water and drive generating turbines; it does not simply burn nuclear fuel.
\item \textbf{Fusion}: nuclei much lighter than iron, such as hydrogen's isotopes deuterium and tritium, combine into heavier ones. They too move toward the valley bottom, on a steeper slope with greater energy release. The Sun is a fusion reactor operating for 4.6 billion years, not a pile of fire: about 600 million tonnes of hydrogen fuse into helium each second. Earthly fusion power remains an engineering challenge seemingly always decades away. Heating nuclei above a hundred million degrees and keeping them confined is much harder than ignition itself.
\end{itemize}

Why are chemical reactions so far behind? Chemistry moves external electrons, involving a few electronvolts per atom; nuclear changes involve millions of electronvolts per nucleus, roughly a millionfold difference. Thus sugar-cube-sized nuclear fuel can replace a truckload of coal, and nuclear weapons have terrifying energy density. This again reminds us that nuclei are beyond chemical changes' reach. Every later discussion of bonds, reactions, and metabolism assumes nuclei quietly remain in place.

\section{Radioactivity, Dose, and Risk: Three Different Things}

Before discussing nuclear matters, separate three frequently confused concepts. Understanding them prevents both panic over radioactive bananas and carelessness about real dangers.

\textbf{Radioactivity} characterizes a substance: how many nuclear decays occur per second. The unit is becquerel, Bq, one decay per second. A gram of radium is about $3.7\times 10^{10}$ Bq, a banana about 15 Bq from potassium-40, and your body a few thousand Bq, mainly from potassium-40 and carbon-14. Activity describes the source, without yet saying anything about you.

\textbf{Radiation dose} describes energy you actually absorb and its effects. The same source behind lead shielding or against your skin produces entirely different situations. Holding an $\alpha$ source, whose rays cannot penetrate outer skin, differs enormously from inhaling it. A common unit is the sievert, Sv, incorporating radiation type and organ sensitivity. For reference: natural background from soil, cosmic rays, food, and radon averages about 2--3 mSv per person yearly; an international flight about 0.1 mSv; a chest X-ray about 0.02--0.1 mSv; a banana about 0.0001 mSv.

\textbf{Risk} translates dose into the probability of health consequences, and the relationship is complicated. At very low doses, risk is indistinguishable amid everyday risks; at high doses, such as more than 1000 mSv in a short time, injury is certain and serious. Between them, estimates rely on population statistics and contain much uncertainty. Radioactivity does not automatically mean danger, just as cars do not automatically mean crashes. What matters is dose and how it reaches the body. Medicine also follows justification and optimization: a radiation examination must offer clear benefit and use as little dose as possible.

Once these are distinct, legitimate radioactive applications lose their mystery:

\begin{itemize}[itemsep=2pt]
\item \textbf{Medical imaging}: tiny amounts of short-lived radionuclides, such as technetium-99m with a 6-hour half-life, enter the body and collect in particular organs. Emitted $\gamma$ rays leave the body for a camera, giving a functional picture. PET uses pairs of $\gamma$ rays from positron annihilation to reveal tumor metabolism.
\item \textbf{Radiotherapy}: rapidly dividing cancer cells repair radiation damage less effectively than normal cells. Precisely focused $\gamma$ beams from multiple directions concentrate dose at the tumor while spreading it across healthy tissue, honing dose into a sharp, accurate knife.
\item \textbf{Dating}: living things continually exchange carbon with the environment, maintaining a constant carbon-14 fraction. Death stops exchange, leaving carbon-14 only to decrease. Measure the fraction remaining and use its 5730-year half-life to read the age. This moved archaeology from guessing dynasties to calculating years and works back roughly fifty thousand years.
\item \textbf{Nuclear energy}: fission plants emit no carbon dioxide and use fuel of extremely high energy density. The costs are managing long-lived radioactive waste and accidents of very low probability but very high consequence. It is an engineering problem and a classic case of society weighing risk.
\end{itemize}

\section{The Model's Boundaries}

This chapter's picture is useful, but three boundaries deserve clarity.

First, half-life is statistical and \textbf{fails for small atom counts}. Saying iodine-131 has an 8-day half-life implicitly assumes countless atoms. With 4 coins left, the next round might eliminate none or all: the average remains, but one result is unpredictable. Saying a substance disappears completely after 80 days is also wrong. In theory, a tail always remains, merely becoming negligible.

Second, stable and radioactive have no sharp boundary. Some nuclei once considered absolutely stable were later found to decay with half-lives on the order of $10^{20}$ years, hundreds of millions of universe ages. Stability is a matter of degree.

Third, protons, neutrons, and two competing forces form a coarse nuclear model, sufficient for decay and energy but not every detail. It cannot explain why certain magic proton or neutron numbers are especially stable, or why nuclei have shapes and vibrate. The nuclear many-body problem remains an active, incompletely solved research frontier.

\textbf{Translator safety note:} The statement below about radiotherapy patients needs qualification: external-beam treatment does not make a patient radioactive, but internal or systemic radioactive treatments may require temporary precautions. Follow the treating team's instructions. See \href{https://www.cancer.gov/about-cancer/treatment/types/radiation-therapy}{National Cancer Institute guidance}.

\begin{wuqu}
``Anything exposed to radiation becomes radioactive itself.'' This widespread misconception makes patients fear radiotherapy and consumers reject irradiated foods. In fact, after $\gamma$ or X-rays strike an object, energy is absorbed and the process ends, leaving no lingering radiation behind. You are not radioactive after an X-ray; radiotherapy patients are not radioactive; neither are irradiated spices or medical equipment. Ordinary material acquires radioactivity only by contamination with radioactive dust or liquid, or by intense neutron bombardment transforming stable nuclei into radioactive ones, activation. The latter requires reactor-level conditions absent from daily life. Distinguish irradiation from contamination: a knife that cuts a radioactive sample may collect radioactive dust and require prescribed cleaning, but a knife merely exposed beside a source stays clean.
\end{wuqu}

\section{Back to the Smoke Alarm and Banana}

\textbf{Translator safety note:} A long radioactive half-life does not make a smoke alarm maintenance-free. Test household alarms monthly and replace them according to the manufacturer's instructions; USFA recommends replacing the entire alarm after ten years. Never dismantle its radioactive source. See \href{https://www.usfa.fema.gov/prevention/home-fires/prepare-for-fire/smoke-alarms/}{USFA smoke-alarm guidance}.

Now for our opening promise. The tiny amount of americium-241 in an ionization smoke alarm, usually under a microgram, continuously emits $\alpha$ particles. These ionize air molecules in a small chamber, producing a steady weak current between electrodes. Smoke particles entering the chamber capture ions, the current falls, and the circuit sounds the alarm. Seeing smoke depends on alpha decay's continuous output. Why americium-241? Its 432-year half-life means barely declining output over decades, allowing years without maintenance. Alpha radiation cannot penetrate the plastic casing, and sleeping beneath it adds less dose than a few extra bananas per year. This is no reason to dismantle it: an exposed source presents a completely different risk. Use specialist recycling channels.

The banana provides reassurance by counterexample. Potassium-40's half-life is about 1.25 billion years; a banana undergoes roughly 15 decays per second and gives about 0.0001 mSv. Your body contains about 140 grams of potassium, whose potassium-40 makes you intrinsically radioactive by several thousand Bq. This has been true since birth and will remain true all your life. Radioactivity is no monster invented by modern industry; it is part of this planet's and your body's original equipment. Dose deserves respect, rather than the word radioactivity alone.

\section{The Nuclear Story Is Not Finished}

This chapter admitted an unsettling fact: nuclei change, elements transmute, and even stability is relative. Yet that fact provides powerful cancer treatments, clocks for buried artifacts, power plants without black smoke, and a key to understanding the universe.

The larger unresolved question is where today's hundred-odd elements originated if nuclei split and combine, decay and form. The Big Bang made only the lightest few. Calcium in bones, iron in blood, and rare earths in phones accumulated through repeated nuclear changes in stellar fusion furnaces and the violent explosions of dying stars. This chapter's nuclear physics is the universe's element-manufacturing line. How it starts, stops, and scatters products into space is the trail Chapter 16, ``Where Elements Come From and Where They Go,'' will follow to the ends of stellar lives.

\begin{tiaozhan}
How much energy does one fission release? A uranium-235 fission yields about 200 MeV, megaelectronvolts. In joules, $200\times 10^{6}\times 1.6\times 10^{-19}\approx 3.2\times 10^{-11}$ J. Tiny, apparently, but 1 gram contains about $6.02\times 10^{23}\div 235\approx 2.6\times 10^{21}$ atoms. Fissioning all releases about $8\times 10^{10}$ J, or about $8\times 10^{13}$ J per kilogram. Good coal releases about $3\times 10^{7}$ J per kilogram; division makes a kilogram of fully fissioned uranium-235 equivalent to about 2700 tonnes of coal. The millionfold gap arises from Einstein's $E=mc^{2}$. Fission fragments' total mass is about a thousandth less than the original nucleus, and that missing mass converts to energy through the astronomical factor $c^2$. Chemical reactions change mass too, but only on the order of a billionth, unmeasurable by a balance.
\end{tiaozhan}

\section*{Try It Yourself}

\begin{sikaoti}
\item Radium-226 becomes radon through one $\alpha$ decay. Write the equation: radium is element 88 and mass number decreases by 4; check the periodic table for the product's atomic number. Why does its chemistry differ completely from radium's?
\item A radionuclide has a 5-day half-life. What fraction remains after 20 days? Sketch remaining amount against time and mark each half-life. Is it valid to say four times five is twenty, so everything has decayed? Where is the mistake?
\item Why need an $\alpha$ source sitting on a table not cause excessive panic, while grinding it to powder and inhaling it is extremely dangerous? Give one reason involving dose and one involving radiation type.
\item Draw a tug-of-war diagram with a small nucleus at left and a very large one at right. Use arrows for strong interactions joining neighbors and electromagnetic repulsion joining all protons. Explain why small nuclei are stable and oversized ones inevitably unstable. Note that real nuclei contain no sticks or ropes; this is a human representation.
\item Someone says nuclear radiation is so frightening that we should live in zero-radiation surroundings. Respond using radioactivity, dose, and risk. Is zero radiation possible on Earth, and is eliminating natural background a worthwhile goal?
\item An excavated skeleton contains one eighth of living organisms' carbon-14 level. Estimate its age and state two assumptions behind the method. (Hint: consider the source of living carbon-14 levels and what happens after death.)
\end{sikaoti}
